\section{研究背景与问题定义}

本章的目的在于说明研究问题的来源和建模边界。移动设备、社交媒体平台、即时通知和在线学习工具已经嵌入日常生活，用户的设备使用时长、手机解锁次数、通知暴露、社交媒体使用、睡眠质量和身体活动之间可能存在复杂关联。对这些变量进行数据分析，有助于从课程实践角度理解数字生活方式数据如何支持风险筛查、行为指标预测和用户画像构建。

需要强调的是，本文使用的数据为合成教学/基准数据，研究目标是完成《大数据分析与应用》课程中的数据处理、可视化和机器学习建模训练。报告中的“高风险”仅对应数据集中的 \texttt{high\_risk\_flag} 标签，不代表医学诊断结论；模型输出也不能直接用于个体健康判断或自动化干预。

\subsection{研究问题}

围绕数字生活方式变量，本文设置四个研究问题：

\begin{itemize}
  \item Q1：在排除心理状态、效率得分和数字依赖得分等潜在泄漏字段后，能否对 \texttt{high\_risk\_flag} 进行高风险数字生活方式筛查？
  \item Q2：使用人口背景、数字行为、生活习惯和工程特征，能否稳定预测 \texttt{digital\_dependence\_score}？
  \item Q3：\texttt{productivity\_score} 是否能被当前可观测特征稳定预测，若不能，负结果本身说明什么？
  \item Q4：只基于数字行为与生活习惯数值特征，能否形成具有解释价值的探索性用户画像？
\end{itemize}

\subsection{任务设计}

为回答上述问题，本文将实验组织为三类机器学习任务。分类任务预测 \texttt{high\_risk\_flag}，重点关注 Recall、F1、PR-AUC 和混淆矩阵，因为高风险筛查中漏判会降低模型的实际参考价值。回归任务以 \texttt{digital\_dependence\_score} 为主线，同时保留 \texttt{productivity\_score} 作为对照，以检验不同目标变量在同一特征体系下的可预测性差异。聚类任务只使用数字行为和生活习惯数值特征，在不使用结果变量参与训练的前提下形成用户画像，再用风险比例、数字依赖得分和效率得分进行事后解释。

\subsection{实验约束}

全部实验使用 Python、pandas、numpy、scipy、scikit-learn 与 matplotlib，在 CPU 环境下运行，不使用深度学习模型。随机过程统一设置 \texttt{random\_state=42}，以保证结果可复现。正文中所有指标均来自 \texttt{results/} 目录下已有 CSV 或文本文件，不对未运行的实验作推断。
